% !TEX root = ../main.tex

\chapter{Conclusion}\label{chapter:conclusion}

This thesis investigated uncertainty quantification for Graph Neural Network surrogates of agent-based transport models on a Paris road network.\ The work used a fixed $1{,}000$-scenario subset of the $10{,}000$-scenario Natterer corpus and evaluated eight base trials, three uncertainty-aware extensions, a five-member Deep Ensemble, six UQ families, and three post-hoc calibration and deployment tools on a 100-graph test set with $3{,}163{,}500$ node-level predictions.\ The findings are specific to the Paris capacity-reduction setup studied here; the architecture (PointNetTransfGAT) and the scenario subset are held fixed throughout.\ The detailed research-question answers appear in Section~\ref{sec:rq_answers}, the primary findings in Section~\ref{sec:primary_results}, and the threats to validity that bound these claims in Section~\ref{sec:threats_to_validity}.

Machine-learning surrogates can collapse a multi-hour MATSim run to seconds of GNN inference, but the speed advantage only translates into trustworthy decision support when the predictions come with a principled statement of reliability.\ The thesis showed that effective uncertainty quantification for GNN traffic surrogates is achievable at a tractable cost.\ MC Dropout supplies a ranking signal that holds up across scenarios; conformal prediction turns that signal into an empirical marginal coverage statement on the held-out evaluation graphs; and a frozen-backbone quantile head delivers calibrated intervals at a single deterministic forward pass.\ Compute costs differ across the toolkit -- MC Dropout multiplies inference by a factor of $S$, the Deep Ensemble multiplies training by roughly $5\times$, and the frozen-backbone CQR head adds only $134$ trainable parameters at deployment -- yet all three are practical inside an offline policy-evaluation loop.\ MC Dropout and conformal prediction sit post-hoc on top of the trained backbone, while the CQR variant additionally trains a small new head.

A few ideas are worth keeping separate here.\ Ranking is not calibration, calibration is not coverage, and a backbone trained for one objective is not always a good starting point for a head trained under another.\ The freezing-backbone observation is a concrete example of that last point: in the experiments here, simply keeping the MSE-trained backbone frozen was enough to flip the CQR result from negative to positive.\ That makes it a sensible starting point for further work on uncertainty-aware extensions for this class of surrogates, with broader architectural and dataset coverage than was feasible inside the present scope.

To contextualise the GNN UQ stack against simpler tabular learners, three non-GNN baselines are reported in Section~\ref{sec:non_gnn_baselines}: Random Forest, XGBoost and a single-hidden-layer MLP, all on the same scenario-level held-out test set as T8 with default-style configurations from each model's official documentation.\ Both tree baselines exceed T8 on point accuracy ($R^2$ of $0.66$ for Random Forest and $0.74$ for XGBoost, against $0.60$ for T8), in line with the tabular-learning literature for sparse-response targets.\ The focus of this thesis is on reliability around each prediction rather than peak point accuracy: a per-prediction uncertainty pipeline -- MC Dropout, regression $\sigma$-scaling, conformal coverage and frozen-backbone CQR -- that the default tree baselines do not provide, because the methods are tightly coupled to architectural dropout, soft probabilistic outputs, and a differentiable backbone with continuous representations.\ Tree-native UQ extensions such as quantile regression forests and NGBoost are the obvious next comparison to run; this work does not claim that no tree-based pipeline can match the GNN UQ stack, only that the default configurations tested here do not.
